% SPDX-FileCopyrightText: Copyright (c) 2020-2026 Yegor Bugayenko
% SPDX-License-Identifier: MIT

\section{Introduction}

Modern object-oriented languages are large and still growing.
C++, Java, C\#, Python, Ruby, and JavaScript each offer hundreds of
  features, and every new release adds more.
The more features a language has, the harder its programs are to reason
  about, to optimize, and to search for bugs, because every feature may
  interact with every other one~\citep{dunsmore1978programming,
  mccabe1976complexity, basili1984software, chambers1989customization}.
Functional programming rests on a small and well-understood foundation,
  the \(\lambda\)-calculus, into which its many features can be lowered
  and then analyzed uniformly.
Object-oriented programming has no such foundation in common use: we know
  of no minimal object model to which an arbitrary object-oriented program
  can be lowered.
\eolang{} was earlier proposed as exactly such a
  foundation~\citep{bugayenko2021eolang}, yet it was never demonstrated
  how programs written in Java, C++, or other mainstream languages map
  onto it.
This paper fills that gap.

The work of \citet{bugayenko2022origins} explains the design of the \eolang{}
  standard library of objects, which are used below.
